% Please do not change %
\documentclass[11pt]{article} % font size
\usepackage[margin=1in]{geometry} % margins
\usepackage{amsmath,amsthm,amssymb} % for the  common "math symbols"
\usepackage[shortlabels]{enumitem} % for enumeration
\usepackage{booktabs} % if you need tables
\usepackage{listings}
\usepackage{color}

\lstset{
  frame=none,
  xleftmargin=2pt,
  stepnumber=1,
  numbers=left,
  numbersep=5pt,
  numberstyle=\ttfamily\tiny\color[gray]{0.3},
  belowcaptionskip=\bigskipamount,
  captionpos=b,
  escapeinside={*'}{'*},
  language=haskell,
  tabsize=2,
  emphstyle={\bf},
  commentstyle=\it,
  stringstyle=\mdseries\rmfamily,
  showspaces=false,
  keywordstyle=\bfseries\rmfamily,
  columns=flexible,
  basicstyle=\small\sffamily,
  showstringspaces=false,
  morecomment=[l]\%,
}
% Feel free to include additional packages (if needed), but make sure it does not override the margin/font requirements.
%
\setlength{\parindent}{0pt} % no indent for new paragraphs.
\setlength{\parskip}{5pt} % distance between paragraphs, which will be used when you have a blank line in your LaTeX code.


%%%%%%%%%%%%%%%%%%%%%%%%%%%%%%%%%%%%%%%%%%%%%%%%%%%%%%%%%%%%%%%%%%%%%%%%%%%%%%
\title{Principles of Programming Languages: \\Problem Sheet 2 - Memory and monads}
%
\author{}
\date{}
%
\begin{document}
%
\maketitle

%You may use the usual \LaTeX\ environments to structure your answers:
%\begin{align}
%        \label{basegnn}
%        \mathbf{h}_u^{(t+1)} = \sigma \bigg( \mathbf{W}_\text{self}^{(t)} \mathbf{h}_u^{(t)}  + \mathbf{W}_\text{neigh}^{(t)} \sum_{v \in N(u)} \mathbf{h}_v^{(t)} + \mathbf{b}^{(t)} \bigg). 
%\end{align}


%\begin{table}[h]
%\caption{A simple table.}
%\centering
%\begin{tabular}[t]{lrrrr}
%\toprule
%Graphs & Can & Be  & Fun & Too \\ 
%\midrule 
%$G_1$ & $1$  & $2$ & $3$ & $4$ \\ 
%$G_2$ & $-1$  & $-2$ & $-3$ & $-4$ \\ 
%$G_3$ & $0.1$  & $0.2$ & $0.3$ & $0.4$ \\ 
%\bottomrule
%\end{tabular}
%\label{tab}
%\end{table}

\section*{Question 1}

% To answer each part of the question, please simply use a paragraph command 
% \paragraph{(a)}
\lstinputlisting[language=Haskell, lastline=15]{../PS2.fun}

\section*{Question 2}

\lstinputlisting[language=Haskell, linerange={17-34}]{../PS2.fun}

\clearpage
\section*{Question 3}
\lstinputlisting[language=Haskell, lastline=11]{../PS2.hs}

\section*{Question 4}
\lstinputlisting[language=Haskell, linerange={13-14}]{../PS2.hs}

\section*{Question 5}
\begin{lstlisting}[language=Haskell]
-- Q5
-- Suppose xm = \mem -> (x, mem), f x = \mem -> (y, mem), g y = \mem -> (z, mem)
(xm $> f) $> g 
= \mem -> let (x, mem') = xm mem in f x mem' $> g 
= \mem -> let (x, mem') = xm mem in let (y, mem'') = f x mem' in g y mem''
xm $> (\x -> f x $> g) 
= xm $> \x -> (\mem -> let (x', mem') = x mem in let (y, mem'') = f x' mem' in g y mem'')
= \mem -> let (x', mem') = xm mem in let (y, mem'') = f x' mem' in g y mem''

result x $> f
= \mem -> (x, mem) $> f
= \mem -> let (y, mem') = x mem in f y mem' 
= \mem -> f x mem
= f x

xm $> result
= \mem -> let (x', mem') = xm mem in result x' mem'
= \mem -> let (x', mem') = (x, mem) in result x' mem'
= \mem -> result x mem
= \mem -> (x, mem)
= xm
\end{lstlisting}

\clearpage
\section*{Question 6}
\lstinputlisting[language=Haskell]{../MonadMemPrint.hs}

\clearpage
\section*{Question 7}
\paragraph{(a)}
\lstinputlisting[language=Haskell]{../EvalMem.fun}

\paragraph{(b)}
Update \texttt{evalargs} to evaluate parameters right to left, 
and \texttt{eval (Apply f es)} to evaluate the function after.
\lstinputlisting[language=Haskell, linerange={16-28}]{../PS2.hs}

\clearpage
\section*{Question 8}
\paragraph{(a)}
\begin{lstlisting}[language=Haskell]
  (e_1; e_2); e_3
  = e_1; ((lambda x -> e_2 x); e_3)
  = e_1; ((lambda x -> e_2 x); e_3; result)
  = e_1; ((lambda x -> e_2 x); (lambda y -> e_3 y; result))
  = e_1; ((lambda x -> e_2 x); (lambda y -> e_3 y))
  = e_1; (lambda x -> e_3 (e_2 x))
  = e_1; (e_2; e_3)
\end{lstlisting}
\begin{lstlisting}[language=Haskell]
  nil; e1
  = (nil; result); e1
  = nil; (lambda x -> result x; e1)
  = nil; (lambda x -> e1 x)
  = e1
\end{lstlisting}
\texttt{nil} may not satisfy the type of the monad.
\begin{lstlisting}[language=Haskell]
  e1; nil
  = (e1; result); nil
  = e1; (lambda x -> result x; nil)
  = e1; (lambda x -> nil x)
  = nil
\end{lstlisting}

\paragraph{(b)}
From Tutorial 1, we know
\begin{lstlisting}[language=Haskell]
let val x = e1 in e2
= (lambda (x) e2)(e1)
\end{lstlisting}
Then it is shown that the lambda expression is equivalent.
\begin{lstlisting}[language=Haskell]
  e1; e2 
  = (e1; e2); result
  = e1; (lambda x -> e2 x; result)
  = e1; lambda x -> e2 x;
  = (lambda x -> e2)(e1)
\end{lstlisting}

\section*{Question 9}
\lstinputlisting[language=Haskell, linerange={30-36}]{../PS2.hs}

\clearpage
\section*{Question 10}
\lstinputlisting[language=Haskell, firstline=38]{../PS2.hs}

\end{document}